% Revize výukové přípravy, výchozí soubor GitHub Work-MAT, 88. hodina.
% Zdroje: Realisticky.cz 020410; grafické dělení definičního oboru z hodin 70–71
% Úlohy jsou vlastní / úlohy výchozí přípravy, případně nově didakticky obměněné.
\documentclass[11pt,a4paper]{article}
\input{../../sablony/pripravy/priprava.sty}
\usepackage{tikz}
\nastavPripravu{1. ročník}{Rovnice a nerovnice}{88}{Rovnice se dvěma absolutními hodnotami}
\begin{document}
\blok{Návaznost a algoritmus}
Při grafu $y=|x-4|+|x+4|$ jsme již určili kritické body $-4$ a $4$. U rovnice postupujeme stejně: seřadíme nulové body vnitřních výrazů, rozdělíme osu, odstraníme moduly a každého kandidáta zkontrolujeme v příslušném intervalu.
\blok{Dvě možné metody}
\begin{zadaniBox}[title={Vlastnost absolutní hodnoty}]
\textbf{Rovnost dvou absolutních hodnot.} Pro reálné výrazy $A,B$ platí
\[
|A|=|B|\quad\Longleftrightarrow\quad A=B\ \text{nebo}\ A=-B.
\]
\end{zadaniBox}
Pro součet $|A(x)|+|B(x)|=c$ obvykle volíme \textbf{metodu kritických bodů}; zkratku z definice nelze přenést na součet absolutních hodnot.
\begin{center}\small
\begin{tabular}{@{}lccc@{}}\toprule
Oblast & $x<-4$ & $-4\le x\le4$ & $x>4$\\\midrule
$|x-4|$ & $4-x$ & $4-x$ & $x-4$\\
$|x+4|$ & $-x-4$ & $x+4$ & $x+4$\\\midrule
Součet & $-2x$ & $8$ & $2x$\\\bottomrule
\end{tabular}
\end{center}
\blok{Společně řešené příklady}
\begin{prikladReseniBox}
Vyřešte $|x-4|+|x+4|=10$.
\tcblower
\[\begin{aligned}
x<-4&:\quad -2x=10\ \Rightarrow\ x=-5 \\
-4\le x\le4&:\quad 8\ne10 \\
x>4&:\quad 2x=10\ \Rightarrow\ x=5 \\
K&=\{-5;5\}
\end{aligned}\]
Zkouška: pro $x=-5$ i $x=5$ vyjde $10$.
\end{prikladReseniBox}
\begin{prikladReseniBox}
Vyřešte $|x-1|=|2x+1|$ úspornější metodou.
\tcblower
\[\begin{aligned}
x-1&=2x+1\ \Rightarrow\ x=-2 \\
x-1&=-(2x+1)\ \Rightarrow\ 3x=0\ \Rightarrow\ x=0 \\
K&=\{-2;0\}
\end{aligned}\]
Oba kořeny ověř dosazením; nevytváříme tři intervaly, když stačí dvě rovnice.
\end{prikladReseniBox}
\begin{prikladReseniBox}
Pro jaká $x$ platí $|x-1|+|x+1|=2$?
\tcblower
\[\begin{aligned}
x<-1&:\quad -2x=2\ \Rightarrow\ x=-1\ \not<-1 \\
-1\le x\le1&:\quad 2=2 \\
x>1&:\quad 2x=2\ \Rightarrow\ x=1\ \not>1 \\
K&=\langle-1;1\rangle
\end{aligned}\]
Výsledkem je celý interval, ne pouze dvě krajní hodnoty.
\end{prikladReseniBox}
\blok{Samostatná práce}
\begin{zadaniReseniBox}
$|x-3|+|x+1|=6$.
\tcblower
\[\begin{aligned}
x<-1&:\quad -2x+2=6\Rightarrow x=-2 \\
-1\le x\le3&:\quad 4\ne6 \\
x>3&:\quad 2x-2=6\Rightarrow x=4 \\
K&=\{-2;4\}
\end{aligned}\]

\end{zadaniReseniBox}
\begin{zadaniReseniBox}
$|x-5|=|x+3|$.
\tcblower
\[\begin{aligned}
x-5&=-(x+3) \\
2x&=2 \\
K&=\{1\}
\end{aligned}\]

\end{zadaniReseniBox}
\begin{zadaniReseniBox}
$|x+2|+|x-2|=4$.
\tcblower
\[\begin{aligned}
-2\le x\le2&:\quad 4=4 \\
x<-2&:\quad -2x=4\Rightarrow x=-2\not<-2 \\
x>2&:\quad 2x=4\Rightarrow x=2\not>2 \\
K&=\langle-2;2\rangle
\end{aligned}\]

\end{zadaniReseniBox}
\blok{Volitelná rezerva}
\begin{zadaniReseniBox}
Najděte řešení $|2x-1|+|x+2|=6$ (kritické body $-2$ a $\frac12$).
\tcblower
\[\begin{aligned}
x<-2&:\quad (-2x+1)+(-x-2)=6\Rightarrow x=-\frac73 \\
-2\le x<\frac12&:\quad (-2x+1)+(x+2)=6\Rightarrow x=-3\notin\langle-2;\frac12) \\
x\ge\frac12&:\quad (2x-1)+(x+2)=6\Rightarrow x=\frac53 \\
K&=\left\{-\frac73;\frac53\right\}
\end{aligned}\]

\end{zadaniReseniBox}
\blok{Domácí úkol}
\begin{zadaniReseniBox}[title={Domácí příprava}]
\begin{enumerate}[label=\alph*),itemsep=2pt]
\item $|x-2|+|x+2|=8$.
\item $|x+1|=|x-5|$.
\item $|2x-1|=5$.
\item $|x-5|+|x+1|=10$.
\item $|x-4|=|x+2|$.
\item $|x-2|+|x+2|=4$.
\end{enumerate}
\tcblower
\begin{enumerate}[label=\alph*),itemsep=1pt]
\item
\[\begin{aligned}
x<-2&:\quad -2x=8\Rightarrow x=-4 \\
-2\le x\le2&:\quad4\ne8 \\
x>2&:\quad2x=8\Rightarrow x=4 \\
K&=\{-4;4\}
\end{aligned}\]
\item
\[\begin{aligned}
x+1&=-(x-5) \\
2x&=4 \\
K&=\{2\}
\end{aligned}\]
\item
\[\begin{aligned}
2x-1&=\pm5 \\
x&=-2\ \text{nebo}\ x=3 \\
K&=\{-2;3\}
\end{aligned}\]
\item
\[\begin{aligned}
x<-1&:\quad -2x+4=10\Rightarrow x=-3 \\
-1\le x\le5&:\quad6\ne10 \\
x>5&:\quad2x-4=10\Rightarrow x=7 \\
K&=\{-3;7\}
\end{aligned}\]
\item
\[\begin{aligned}
x-4&=-(x+2) \\
2x&=2 \\
K&=\{1\}
\end{aligned}\]
\item
\[\begin{aligned}
x<-2&:\quad -2x=4\Rightarrow x=-2\not<-2 \\
-2\le x\le2&:\quad4=4 \\
x>2&:\quad2x=4\Rightarrow x=2\not>2 \\
K&=\langle-2;2\rangle
\end{aligned}\]
\end{enumerate}
\end{zadaniReseniBox}
\poznamka{45 minut: kritické body 6, první vzor 12, rovnost dvou modulů 7, interval řešení 6, samostatná práce 10, reflexe 4. Náročná rezerva a poslední DÚ jsou dobrovolné.}
\end{document}
